\chapter{Implementation Details}
\label{ch:implementation}

\begin{center}
\textit{This chapter details the implementation: the technology stack, the
repository structure, the NestJS backend, the Python workers, the OCR service,
the web client, and the deployment infrastructure.}
\end{center}

\vspace*{\fill}
\clearpage
\section{Technology Stack}

\noindent Table~\ref{tab:stack} summarises the implemented technology stack.

\begin{table}[htbp]
\centering
\small
\begin{tabular}{@{}p{4.3cm}p{10.2cm}@{}}
\toprule
\textbf{Component} & \textbf{Technology} \\
\midrule
API framework & NestJS (TypeScript, ES2022, NodeNext modules, strict mode) \\
Web client & Next.js (App Router) \\
OCR service & FastAPI (Python 3.12+) \\
Workers & Python 3.12+, pika (direct AMQP consumption) \\
Database & PostgreSQL 17, Drizzle ORM \\
Message broker & RabbitMQ 3 \\
AI integration & Internal OpenAI-compatible gateway (OmniRoute) \\
Validation & Zod (TypeScript), Pydantic (Python) \\
OCR engines & PyMuPDF (embedded text), PaddleOCR (fallback) \\
Deployment & Docker Compose, nginx, Cloudflare Tunnel \\
Build & pnpm workspaces, Turborepo \\
\bottomrule
\end{tabular}
\caption{Implemented technology stack.}
\label{tab:stack}
\end{table}

\section{Repository Structure}

\noindent The monorepo is organised as follows:

\begin{itemize}
  \item \texttt{apps/api} --- the NestJS modular monolith (twenty feature
  modules plus application, database, and health scaffolding);
  \item \texttt{apps/web} --- the Next.js client;
  \item \texttt{apps/ocr} --- the FastAPI OCR service;
  \item \texttt{apps/workers} --- the Python async workers
  (\texttt{worker/} AI and material consumers, \texttt{ocr-worker/} pull
  client);
  \item \texttt{packages/contracts} --- shared API contracts and Zod schemas,
  imported by both web and API;
  \item \texttt{packages/database} --- Drizzle schema, 48 migrations, and the
  database client;
  \item \texttt{packages/auth}, \texttt{packages/ai}, \texttt{packages/shared}
  --- shared authentication, AI, and common utilities;
  \item \texttt{infrastructure/compose} --- the Docker Compose definitions;
  \item \texttt{docs} --- architecture, proposal, status, tracker, and
  user-validation documentation.
\end{itemize}

\noindent TypeScript is compile-strict: ESM with explicit \texttt{.js} import
extensions, no \texttt{any} without an explicit exemption, and a single
\texttt{pnpm typecheck} covering all workspaces. Shared contracts keep the web
client and the API on the same DTOs, while Pydantic models on the Python side
re-validate worker inputs at the boundary.

\section{Backend (NestJS API)}

\noindent Each feature area is a standalone NestJS module
(\texttt{*Module}, \texttt{*Controller}, \texttt{*Service}, \texttt{*DTO}).
Cross-cutting concerns are centralised: the global guard and exception
infrastructure (Chapter~\ref{ch:security}), a DTO validation layer built on
Zod schemas from the contracts package, and Drizzle repositories that always
apply an \texttt{institute\_id} filter. Long-running operations from any
controller call \texttt{JobsService}, which provides the insert-then-publish
job lifecycle, cancellation, retry with deduplication, and status queries.
Routing uses one global prefix \texttt{/api/v1}; the health endpoint lives at
\texttt{GET /api/v1/health}.

\section{AI and Material Workers (Python)}

\noindent The Python workers are plain long-running processes that consume
RabbitMQ queues directly over AMQP (pika), rather than a task framework: the
AI consumer binds \texttt{ai\_generation} with bounded concurrency (default 5)
and prefetch, and the material consumer binds \texttt{jobs}. On start, a
worker recovers stale queued messages so that no work is lost across worker
restarts. Generation work resolves READY source materials from the database,
builds a context budget (a bounded character window divided into overlapping
chunks), calls the model through OmniRoute (\texttt{POST
/v1/chat/completions}), parses and Pydantic-validates each response, then
writes results directly to PostgreSQL with provenance and deterministic
de-duplication. Transient failures and transport errors are retried with
backoff (maximum three attempts); terminal errors settle the job as failed.
The material-enhancement pass runs entirely inside the \emph{API} coordinator
(deterministic, no model), not in a worker.

\section{OCR Service and OCR Worker}

\noindent The OCR service is a dependency-lean FastAPI application with no
business-domain knowledge. It exposes \texttt{GET /health} and extraction
logic that, per page, first attempts PyMuPDF embedded text and falls back to
rasterisation + PaddleOCR when the embedded text is too short; its output is
then passed through a deterministic character-normalization pass. The
standalone \texttt{ocr-worker} package is a pull client: it heartbeats, claims
chunks from the API's worker endpoints, downloads source bytes from the
corresponding read route, runs the extraction range locally, and submits page
results with a page-count field. Legacy one-shot extraction (a file posted to
the OCR service and text returned synchronously into a material) is retained
for compatibility but is not the primary path.

\section{Web Client (Next.js)}

\noindent The web client is built on the App Router and consumes the shared
contracts package. Tenancy is client-side state: the active institute id is
kept in local storage, the API client attaches \texttt{x-institute-id} and the
CSRF header on every request, and a single-flight refresh handles session
renewal with an outcome that never logs the user out on an unavailable
network. UI components are minimal and dependency-light (a small set of
presentational components rather than a component library), and the
institute-switcher shows the resolved institute-domain permissions per
membership so navigation adapts to what the user may do.

\section{Infrastructure and Deployment}

\noindent The single \texttt{docker-compose.yml} defines 12 services:
PostgreSQL~17, Redis~7, RabbitMQ~3, a one-shot migrate runner, the API, the
web client, nginx, the Cloudflare Tunnel sidecar, OmniRoute, the OCR service,
and the two worker images. In the production shape nothing publishes a host
port; a development override binds 127.0.0.1 loopback ports (nginx on
\texttt{8080} by default) for local tooling, and a demo overlay adds seed data
and a worker-to-gateway wiring. Dockerfiles are dependency-first: pnpm
workspace installs run on a shared cache mount so that build layers for
third-party dependencies only rebuild when their manifests change, keeping
whole-stack rebuilds fast while the running containers never hot-reload source
code (this is why every code change to the API, web, or workers requires a
deliberate rebuild and verification step). The nginx configuration splits
paths exactly as described in Chapter~\ref{ch:architecture}; a Cloudflare
Tunnel with one public hostname is the only ingress.